\chapter{Phương pháp tọa độ trong không gian}

\section{Phương trình mặt phẳng}

\exercise Trong không gian $Oxyz$, cho mặt phẳng $(P): 2x - 3y + z - 5 = 0$. Hãy chỉ ra một vectơ pháp tuyến $\vec{n}$ và tìm tọa độ một điểm thuộc mặt phẳng $(P)$.

\exercise Trong không gian $Oxyz$, viết phương trình tổng quát của mặt phẳng $(P)$ đi qua điểm $A(1; -2; 3)$ và vuông góc với đường thẳng nối hai điểm $B(2; 1; -1)$, $C(0; 3; 1)$.

\exercise Trong không gian $Oxyz$, viết phương trình mặt phẳng $(\alpha)$ đi qua hai điểm $M(1; 0; 2)$, $N(-1; 2; 1)$ và vuông góc với mặt phẳng $(\beta): 2x - y + z - 3 = 0$.

\section{Phương trình đường thẳng trong không gian}

\exercise Trong không gian $Oxyz$, cho đường thẳng $d$ có phương trình tham số:
$$\begin{cases} x = 1 + 2t \\ y = -3 - t \\ z = 4 + 3t \end{cases}.$$
Tìm một vectơ chỉ phương $\vec{u}$ của $d$ và chỉ ra một điểm $M$ thuộc đường thẳng $d$.

\exercise Trong không gian $Oxyz$, viết phương trình chính tắc của đường thẳng $\Delta$ đi qua hai điểm $A(2; -1; 1)$ và $B(0; 3; 2)$.

\exercise Trong không gian $Oxyz$, xét vị trí tương đối giữa hai đường thẳng:
\begin{center}
   $d_1: \frac{x - 1}{2} = \frac{y + 1}{1} = \frac{z}{3}$ và $d_2: \begin{cases} x = 1 + t \\ y = -2 + 2t \\ z = 1 - t \end{cases}$.
\end{center}

\section{Công thức tính góc trong không gian}

\exercise Trong không gian $Oxyz$, cho hai đường thẳng $\Delta_1, \Delta_2$ có các vectơ chỉ phương lần lượt là $\vec{u}_1 = (1; 2; 2)$ và $\vec{u}_2 = (-2; 1; 2)$. Tính $\cos(\Delta_1, \Delta_2)$.

\exercise Trong không gian $Oxyz$, tính góc giữa đường thẳng $d: \frac{x - 1}{2} = \frac{y + 1}{-1} = \frac{z}{1}$ và mặt phẳng $(P): x + y + 2z - 5 = 0$.

\exercise Trong không gian $Oxyz$, cho hai mặt phẳng $(P): x + 2y - z + 1 = 0$ và $(Q): 2x - y + z + 3 = 0$. Tính góc giữa hai mặt phẳng $(P)$ và $(Q)$.

\exercise Trong không gian $Oxyz$, cho ba điểm $A(2; 0; 0)$, $B(0; 3; 0)$, $C(0; 0; 4)$ và điểm $D(2; 3; 4)$. Tính góc giữa đường thẳng $AD$ và mặt phẳng $(ABC)$.

\section{Phương trình mặt cầu}

\exercise Trong không gian $Oxyz$, xác định tọa độ tâm $I$ và bán kính $R$ của mặt cầu $(S)$ có phương trình:
$$(x - 2)^2 + (y + 1)^2 + z^2 = 16.$$

\exercise Trong không gian $Oxyz$, viết phương trình mặt cầu $(S)$ có đường kính $AB$ với $A(1; -2; 3)$ và $B(3; 0; -1)$.

\exercise Trong không gian $Oxyz$, viết phương trình mặt cầu $(S)$ có tâm $I(1; 2; -1)$ và tiếp xúc với mặt phẳng $(P): 2x - 2y + z - 6 = 0$.